\section{总结与延伸}

\subsection{全讲结论汇总}
\begin{longtable}{p{3.0cm}p{3.8cm}p{4.4cm}}
\toprule
\textbf{问题} & \textbf{课程结论} & \textbf{实践含义} \\
\midrule
为什么分数会波动？ & 分数本质是随机变量，噪声可预测但不可忽略 & 必须报告方差与置信区间 \\
噪声来自哪里？ & 数据层、协议层、评审层共同作用 & 先分层诊断，再做治理 \\
如何比较模型？ & 先保证 comparability，再谈 superiority & 不同协议下分数不可直接排序 \\
Agent 评测有什么不同？ & 轨迹级方差与环境噪声显著更高 & 需要环境版本化与可靠性指标 \\
如何在团队落地？ & 建立抗噪声评测 SOP & 从重复评测和协议冻结开始 \\
\bottomrule
\end{longtable}

\subsection{一句话复盘}
\begin{importantbox}{Takeaway}
如果不对 noise 建模，benchmark 会奖励偶然性；对 noise 建模后，benchmark 才会奖励真实进展。
\end{importantbox}

\subsection{拓展阅读}
\begin{itemize}
    \item Percy Liang et al., HELM（Holistic Evaluation of Language Models）相关论文与报告
    \item ``Beyond the Imitation Game''（LM 评测框架讨论）
    \item 统计学习中的 bootstrap 与 uncertainty quantification 经典教材
    \item 针对 LLM-as-judge 的 bias / consistency 研究工作
    \item Agent benchmark（WebArena / SWE-bench / BrowserGym 类）中的复现与评测协议文档
\end{itemize}

\end{document}
